\chapter{问题重述与分析}

\section{问题背景}

边坡失稳通常经历缓慢蠕变、加速变形和快速滑移三个阶段。实际监测中，传感器误差、降雨入渗、孔隙水压力变化、微震活动以及爆破扰动会共同作用于位移曲线，使得“识别阶段、解释扰动、提前预警”不能只依赖单一指标。C题给出了5组相关数据，要求完成校准、分段、异常处理、预测和预警变量筛选。

本文的目标不是单纯追求某一模型的训练精度，而是形成一条能被复核的证据链：数据清洗有依据，阶段节点有多源证据，预测曲线符合工程规律，预警阈值具有持续时间和提前量解释。

\section{五个问题的建模要求}

问题一要求将两类位移传感器数据统一到同一量测基准。由于传感器响应大体线性，但局部点可能受噪声影响，校准模型应稳健且不宜过度弯曲。

问题二要求识别边坡变形三阶段。累计位移具有单调趋势，但阶段转换更多体现为速度水平和加速度趋势的变化，因此需要区分“趋势开始改变”和“速度显著跃迁”。

问题三要求处理多源监测数据的缺失和异常，并分析各变量对表面位移的影响。降雨量、微震事件数是稀疏驱动量，孔压和位移是连续状态量，二者不能使用完全相同的异常规则。

问题四要求利用训练集规律预测实验集指定时刻的表面位移。题目给出实验集阶段标签，这说明阶段内相对演化位置比绝对日历时间更重要，预测模型必须保持阶段规律一致和位移整体单调。

问题五要求从6个候选变量中选择5个并建立预警机制。预警不能只给出一个分位数阈值，还应说明触发持续时间、多源扰动条件和快速阶段提前量。

\section{总体建模思路}

本文采用“鲁棒校准--复合证据分段--变量专属异常识别--阶段归一化残差预测--速度预警闭环”的路线：

\begin{enumerate}
  \item 用Huber稳健仿射模型解决传感器尺度统一问题，并用多模型交叉验证确认一次模型足够稳定。
  \item 用Hampel滤波和滚动中位速度识别速度跃迁，再引入位移趋势和持续加速度作为提前证据。
  \item 对连续状态量使用状态空间补齐，对稀疏驱动量保留真实事件强度，避免把降雨和微震误删为坏点。
  \item 对实验集先生成阶段归一化基线，再用扰动残差模型修正局部偏差，保证预测曲线符合边坡变形规律。
  \item 对预警变量进行全组合验证，并以6 h平滑速度、逆速度趋势和扰动增强共同构成三级预警机制。
\end{enumerate}

这一结构的核心优势是分工明确：物理规律由阶段基线和速度阈值承担，机器学习模型只承担残差修正和贡献排序，避免黑箱模型在竞赛论文中难以解释的问题。

